\documentclass[12pt]{article}
\usepackage[utf8]{inputenc}
\usepackage{amsmath}
\usepackage{amssymb}
\usepackage{geometry}
\geometry{a4paper, margin=2.5cm}

\title{Derivasi Formal Hamiltonian Ising dari Teori Utilitas von Neumann--Morgenstern dalam Sistem Finansial}
\author{}
\date{}

\begin{document}

\maketitle

\section{Aksioma von Neumann--Morgenstern dan Teori Utilitas Harapan}

Teori utilitas yang dikembangkan oleh John von Neumann dan Oskar Morgenstern menetapkan fondasi bagi perilaku agen rasional dalam menghadapi ketidakpastian melalui serangkaian aksioma formal. Dalam kerangka ini, preferensi individu terhadap berbagai *lottery* atau prospek peluang dapat direpresentasikan secara matematis melalui fungsi *expected utility*. Sebuah lotere $L$ didefinisikan sebagai himpunan pasangan terurut yang menghubungkan hasil (*outcome*) dengan probabilitas kemunculannya, $L = \{(x_k, p_k)\}_{k=1}^n$.

Secara formal, utilitas harapan $U(L)$ dari sebuah lotere $L$ didefinisikan sebagai jumlah tertimbang dari utilitas setiap *outcome* yang mungkin terjadi. Jika $p_k \in [0, 1] \subset \mathbb{R}$ merepresentasikan probabilitas terjadinya *outcome* $x_k \in \mathbb{R}$, maka fungsi utilitas dapat diekspresikan sebagai berikut:
\begin{equation}
U(L) = \sum_{k=1}^n p_k u(x_k) \tag{1}
\end{equation}
Di sini, $u: \mathbb{R} \to \mathbb{R}$ merupakan fungsi utilitas bernilai riil yang memetakan perolehan ekonomi ke dalam skala kepuasan subjektif, dengan syarat normalisasi probabilitas $\sum_k p_k = 1$. Dalam sistem ekonomi diskrit, seperti keputusan investasi biner, strategi setiap pemain $i$ dapat direpresentasikan oleh variabel biner $s_i \in \{0, 1\}$. Untuk kepentingan pemetaan ke model fisika, variabel ini seringkali ditransformasikan ke dalam representasi *spin* $\sigma_i \in \{-1, 1\}$, di mana hubungan antara keduanya diberikan oleh $\sigma_i = 2s_i - 1$. Transformasi ini menjadi krusial dalam memetakan perilaku agen menuju formalisme mekanika statistik.

\section{Ekspansi Utilitas Interaktif dan Pemetaan Hamiltonian}

Dalam sistem dengan banyak agen yang saling berinteraksi, utilitas seorang pemain $i$ merupakan fungsi multivariat dari konfigurasi strategi seluruh pemain, $U_i(s_1, s_2, \dots, s_N)$. Untuk memahami dinamika interaksi tersebut, fungsi utilitas dapat diekspansi menggunakan deret Taylor diskrit di sekitar titik referensi. Ekspansi ini memungkinkan kita untuk mengurai kontribusi utilitas menjadi komponen individual, interaksi berpasangan, dan gangguan orde tinggi.

Secara formal, ekspansi utilitas pemain $i$ terhadap seluruh variabel strategi $\mathbf{s}$ hingga orde ketiga dapat dinyatakan sebagai:
\begin{equation}
U_i(\mathbf{s}) = U_i(\mathbf{0}) + \sum_j \frac{\partial U_i}{\partial s_j} s_j + \frac{1}{2} \sum_{j,k} \frac{\partial^2 U_i}{\partial s_j \partial s_k} s_j s_k + O(s^3) \tag{2}
\end{equation}
Untuk memetakan sistem ini ke dalam Hamiltonian tunggal, kita mengasumsikan bahwa interaksi tersebut merupakan *potential game*. Dalam kondisi ini, terdapat fungsi potensial global $\Phi(\mathbf{s})$ yang memenuhi korelasi insentif $\Delta U_i = \Delta \Phi$ untuk setiap perubahan strategi individu. Dengan melakukan ekspansi Taylor pada fungsi potensial global di sekitar $\mathbf{s}=\mathbf{0}$, kita memperoleh:
\begin{equation}
\Phi(\mathbf{s}) \approx \Phi(\mathbf{0}) + \sum_i \left( \frac{\partial \Phi}{\partial s_i} \right) s_i + \frac{1}{2} \sum_i \left( \frac{\partial^2 \Phi}{\partial s_i^2} \right) s_i^2 + \sum_{i<j} \left( \frac{\partial^2 \Phi}{\partial s_i \partial s_j} \right) s_i s_j \notag
\end{equation}
Mengingat variabel strategi bersifat biner ($s_i \in \{0, 1\}$), berlaku identitas $s_i^2 = s_i$. Suku konstan $\Phi(\mathbf{0})$ dapat diabaikan karena tidak mempengaruhi titik ekuilibrium optimasi. Dengan demikian, kita dapat mengonsolidasikan seluruh suku linear ke dalam parameter $h_i$ dan suku interaksi ke dalam $J_{ij}$ sebagai berikut:
\begin{equation}
h_i = \left. \frac{\partial \Phi}{\partial s_i} \right|_{\mathbf{0}} + \left. \frac{1}{2} \frac{\partial^2 \Phi}{\partial s_i^2} \right|_{\mathbf{0}}, \quad J_{ij} = \left. \frac{\partial^2 \Phi}{\partial s_i \partial s_j} \right|_{\mathbf{0}} \notag
\end{equation}
Melalui sintesis parameter tersebut, utilitas total sistem $U_{tot}(\mathbf{s}) \equiv \Phi(\mathbf{s})$ dapat dituliskan dalam bentuk kuadratik standar:
\begin{equation}
U_{tot}(s) = \sum_i h_i s_i + \sum_{i<j} J_{ij} s_i s_j \tag{3}
\end{equation}
Dalam perspektif *statistical physics*, maksimisasi utilitas ekuivalen dengan minimisasi energi sistem, di mana Hamiltonian $H$ didefinisikan sebagai negatif dari fungsi utilitas ($H = -U$). Dengan menggunakan representasi *spin* $\sigma_i$, kita memperoleh bentuk standar Hamiltonian Ising:
\begin{equation}
H = - \sum_i h_i \sigma_i - \sum_{i<j} J_{ij}\sigma_i\sigma_j \tag{4}
\end{equation}
Parameter $h_i$ diinterpretasikan sebagai *bias payoff* individu atau pengaruh eksternal, sedangkan $J_{ij}$ merepresentasikan kekuatan interaksi atau korelasi antar strategi pemain.

\section{Integrasi Kendala Global melalui Penalti Lagrangian}

Pada banyak kasus nyata di pasar finansial, optimasi utilitas seringkali dibatasi oleh kendala global, seperti batasan likuiditas atau kuota jumlah pemain yang memilih aksi tertentu. Kendala ini, misalnya $\sum_i s_i = K$, harus diintegrasikan ke dalam fungsi objektif agar solusi yang dihasilkan tetap berada dalam ruang solusi yang valid. Pendekatan klasik untuk menangani masalah ini adalah dengan menggunakan metode *Lagrangian multiplier* yang menyisipkan kendala ke dalam fungsi utilitas utama.

Untuk menjaga struktur kuadratik yang kompatibel dengan sistem Ising atau *Quadratic Unconstrained Binary Optimization* (QUBO), kendala tersebut seringkali direpresentasikan dalam bentuk penalti kuadrat. Fungsi energi baru yang mencakup *penalty term* dengan koefisien penalti $\lambda$ dapat dirumuskan sebagai berikut:
\begin{equation}
H = - \sum_i h_i s_i - \sum_{i<j} J_{ij} s_i s_j + \lambda \left( \sum_i s_i - K \right)^2 \tag{5}
\end{equation}
Penambahan suku kuadratik ini memastikan bahwa setiap penyimpangan dari kendala global akan meningkatkan energi sistem secara signifikan. Dengan demikian, konfigurasi yang meminimalkan Hamiltonian tersebut secara otomatis akan cenderung memenuhi batasan sistem yang telah ditetapkan.

\section{Derivasi Parameter Hamiltonian dari Distribusi Probabilitas}

Dalam analisis data saham, parameter fisik $h_i$ dan $J_{ij}$ dapat diturunkan secara langsung dari distribusi probabilitas bersama (*joint probability distribution*) tanpa memerlukan asumsi model Markowitz. Jika kita memiliki dua aset dengan probabilitas kejadian bersama yang dinyatakan dalam tabel kontingensi (misalnya $p, q, r, s$ untuk kombinasi naik/turun), maka utilitas harapan dapat dipetakan ke dalam bentuk biner. Proses ini dikenal sebagai *inverse Ising problem*, di mana struktur interaksi diekstraksi dari data observasi.

Parameter *bias* individu $h$ dan koefisien kopling interaksi $J$ dapat dihitung melalui hubungan *log-odds ratio* dari probabilitas yang tersedia. Secara matematis, parameter tersebut diberikan oleh:
\begin{equation}
h_A = \frac{1}{2}\ln\left(\frac{p+q}{r+s}\right) \tag{6}
\end{equation}
\begin{equation}
J_{AB} = \frac{1}{4}\ln\left(\frac{ps}{qr}\right) \tag{7}
\end{equation}
Hasil ini menunjukkan bahwa kopling interaksi $J_{AB}$ memiliki kaitan erat dengan metrik ketergantungan statistik. Penurunan ini membuktikan bahwa struktur Hamiltonian Ising muncul secara alami dari prinsip-prinsip dasar teori permainan dan utilitas VNM, memberikan landasan teoritis yang kuat untuk aplikasi *quantum optimization* dalam bidang finansial.

\end{document}
